% Awesome Source CV LaTeX Template
%
% This template has been downloaded from:
% https://github.com/darwiin/awesome-neue-latex-cv
%
% Author:
% Christophe Roger
% Contenu adapté pour : Samuel SIKATI KENMOGNE
% Version : Analyste Fonctionnel SI & Recette
%
% Template license:
% CC BY-SA 4.0 (https://creativecommons.org/licenses/by-sa/4.0/)

% Commandes de mise en forme définies dans yaac-another-awesome-cv.cls :
%   \company{<url>}{<nom>}   entreprise avec lien (petites capitales, pointillé)
%   \companynolink{<nom>}    entreprise sans site officiel
%   \projet{<intitulé>}      intitulé de projet / contexte (poids medium)

%Section: Work Experience
\sectionTitle{Expériences Professionnelles}{\faSuitcase}
\begin{experiences}
  \experience[CDD]
    {Aujourd'hui}   {Analyste Fonctionnel SI \& Assurance Qualité}{\company{https://www.orange.cm}{Orange Cameroun}}{Douala}
    {Nov. 2025} {
                      \projet{Migration d'une plateforme mobile money : analyse fonctionnelle et homologation}
                      \begin{itemize}
                        \item Recueil des besoins auprès des métiers et contribution aux spécifications fonctionnelles
                        \item Formalisation des règles de gestion et des exigences fonctionnelles de la plateforme
                        \item Mise en place de la traçabilité entre exigences et cas de test pour piloter la couverture fonctionnelle
                        \item Validation fonctionnelle avec les métiers en recette utilisateur (UAT), qualification de \textbf{262 anomalies}
                        \item Conception de \textbf{1\,547 cas de test} répartis en chantiers (fonctionnel, E2E, API, rôles, performance)
                        \item Exécution des tests de migration, PRA et performance (Locust), automatisation de la régression en CI/CD
                      \end{itemize}
                      \projet{Mise en place d'une solution de gestion du courrier entrant}
                      \begin{itemize}
                        \item Analyse et cartographie (BPMN) du circuit existant, de la réception du courrier jusqu'à la réponse
                        \item Recueil des exigences métier et paramétrage des workflows de visa et de diffusion, profils et droits
                        \item Réalisation de la recette fonctionnelle (tests d'acceptation) avant le déploiement
                        \item Rédaction des guides utilisateurs et modes opératoires par profil, formation et support au démarrage
                      \end{itemize}
                    }
                    {Besoins métiers,Spécifications,Règles de gestion,Traçabilité,UAT,Mobile Money,Maarch Courrier}
  \emptySeparator
  \experience[Stage]
    {Juin 2025}   {Analyste Concepteur Junior}{\company{https://societegenerale.cm/fr/}{Société Générale Cameroun}}{Douala}
    {Jan. 2025} {
                      \projet{Conception d'une solution de suivi des anomalies opérationnelles bancaires}
                      \begin{itemize}
                        \item Cartographie du processus de traitement des anomalies (BPMN, UML), jusqu'aux plans d'action correctifs
                        \item Conception du modèle de données associé (échéances, relances automatiques, preuves de correction)
                        \item Mise en place de tableaux de bord KPI et automatisation du reporting, jusque-là manuel
                        \item Développement de la solution (.NET Core) et tests de bout en bout avant la mise en production
                        % \item \textbf{Réduction de 25\,\% du temps de traitement} des anomalies et \textbf{de 35\,\% des erreurs manuelles}
                      \end{itemize}
                    }
                    {Cartographie de processus,BPMN,UML,Modèle de données,KPI,Reporting,Workflow,.NET Core,Banque}
  \emptySeparator
  \experience[Stage]
    {Fév. 2024}   {Analyste Fonctionnel Junior}{\company{https://www.boissonsducameroun.com}{Boissons du Cameroun}}{Douala}
    {Sept. 2023} {
                      \projet{Digitalisation du processus d'analyse des causes profondes des incidents}
                      \begin{itemize}
                        \item Recueil des besoins auprès de l'équipe HSE et rédaction des spécifications fonctionnelles
                        \item Analyse du processus d'investigation et modélisation UML/BPMN du workflow d'enquête
                        \item Structuration de la matrice SCAT (DNV) en contacts, causes immédiates et fondamentales, et actions
                        \item Développement d'une application web (Symfony, Flask) dédiée aux investigations d'incidents
                        % [Résultat masqué] \item \textbf{Réduction de 40\,\% du temps d'investigation} et amélioration de la qualité des rapports d'analyse                                              
                      \end{itemize}
                    }
                    {Recueil des besoins,Spécifications fonctionnelles,Analyse de processus,BPMN,UML,SCAT -- DNV,HSE,Agroalimentaire}
  \emptySeparator
  \experience[Stage]
    {Avr. 2023}   {Analyste Fonctionnel QHSE}{\company{https://www.perenco.com}{Perenco Cameroun}}{Douala}
    {Jan. 2023} {
                      \projet{Optimisation du contrôle qualité QHSE par une application mobile et des macros Excel}
                      \begin{itemize}
                        \item Recueil des besoins auprès des équipes QHSE et analyse du processus d'audit existant
                        \item Modélisation UML de l'application de contrôle (cas d'utilisation, activités)
                        \item Conception du modèle de données des contrôles, des non-conformités et des actions correctives
                        \item Développement d'une application mobile Android (Kotlin) et automatisation des contrôles en VBA Excel
                        % \item \textbf{Amélioration de 20\,\% de l'efficacité des audits} et meilleure conformité réglementaire
                      \end{itemize}
                    }
                    {Analyse de processus,Recueil des besoins,QHSE,Audit,Conformité réglementaire,Oil \& Gas,Kotlin,VBA Excel}
  \emptySeparator
  \experience[Stage]
    {Oct. 2022}   {Analyste Concepteur Back-End}{\companynolink{WinMoney Cameroun}}{Yaoundé}
    {Juin 2022} {
                      \projet{Développement de l'API RESTful d'une plateforme de création de contenus visuels}
                      \begin{itemize}
                        \item Conception du modèle de données, des endpoints REST et de la logique métier (Spring Boot)
                        \item Documentation de l'API et tests fonctionnels des endpoints (Postman)
                        % \item Support de \textbf{jusqu'à 500 requêtes/minute} avec un temps de réponse moyen de 300\,ms
                      \end{itemize}
                    }
                    {Modèle de données,API REST,Logique métier,Documentation,Tests d'API,Postman,Spring Boot}
  \emptySeparator
  \experience[Stage]
    {Oct. 2021}   {Analyste Développeur Junior}{\company{https://kiama.cm}{Kiama S.A}}{Yaoundé}
    {Juin 2021} {
                      \projet{Mise en place d'un système de recommandation multicanal (email/SMS) pour \company{https://ekiosque.cm}{ekiosque.cm}}
                      \begin{itemize}
                        \item Analyse du parcours lecteur et définition des règles de recommandation et de ciblage (email, SMS)
                        \item Personnalisation des recommandations par machine learning (PHP-ML) et recette des envois
                        % \item \textbf{Augmentation de 15\,\% de l'engagement utilisateur} et amélioration du taux de conversion
                      \end{itemize}
                    }
                    {Parcours client,Règles de gestion,Multicanal,Email/SMS,Recette,Engagement utilisateur,Machine Learning}
\end{experiences}
 